% GENERATED FILE. Edit canonical lessons, then run npm run generate:books.
% canonical-source-hash: fnv1a64:8f338d97179e0827
% canonical-lessons: IT-C454-salumeria, IT-C454-pescheria, IT-C454-fruttivendolo, IT-C454-pasticceria, IT-C454-rosticceria

\chapter{Il cibo in negozio}
\label{ch:it-il-cibo-in-negozio}
\hlchaptermodality{454}

\begin{chapteropening}
\textbf{By the end of this chapter:} \emph{I can shop for food: the delicatessen, the fishmonger, the greengrocer, the cake shop and the rotisserie.}

Complete the last lesson of chapter 454: I can shop for food: the delicatessen, the fishmonger, the greengrocer, the cake shop and the rotisserie.
\end{chapteropening}

\section[la salumeria]{la salumeria — the delicatessen (for cured meats)}

\label{lesson:IT-C454-salumeria}



\begin{quote}
Before the new one: say the Italian for the herbalist's, then the Italian for the jeweller's.
\end{quote}

\subsection*{You'll want to know: la salumeria}
\textbf{la salumeria} — \textquotedblleft{}the delicatessen (for cured meats)\textquotedblright{}.

\begin{quote}
In salumeria compro prosciutto e formaggio.
\end{quote}

\begin{quote}
\emph{At the deli I buy ham and cheese.}
\end{quote}

\emph{I salumi} — cured meats.

\begin{cousinweb}
\emph{salumeria} is built on \emph{salume} 'cured meat', from \emph{sale} 'salt', which is inherited from Latin \emph{sāl} (EtymDB).
\end{cousinweb}

\begin{grammarlens}[title={What you've built}]
One more word for telling, explaining and describing at B1.
\end{grammarlens}

\subsection*{Your turn}
\begin{itemize}
\raggedright
  \item \emph{Say it:} \emph{la salumeria}
  \item \emph{Say it:} \emph{la salumeria}, once more
  \item \emph{Say it:} a sentence of your own with \emph{la salumeria}
  \item \emph{Recall:} say the Italian for the herbalist's, then the Italian for the jeweller's, then say \emph{la salumeria} again
\end{itemize}

\begin{tcolorbox}[breakable,colback=teal!4,colframe=teal!35!black,title={Before you move on}]
What does la salumeria mean? (\textquotedblleft{}The delicatessen (for cured meats)\textquotedblright{}.) Say it once more.
\end{tcolorbox}
\section[la pescheria]{la pescheria — the fishmonger's}

\label{lesson:IT-C454-pescheria}



\begin{quote}
Before the new one: say the Italian for the jeweller's, then the Italian for the delicatessen.
\end{quote}

\subsection*{You'll want to know: la pescheria}
\textbf{la pescheria} — \textquotedblleft{}the fishmonger's\textquotedblright{}.

\begin{quote}
La pescheria apre presto la mattina.
\end{quote}

\begin{quote}
\emph{The fishmonger's opens early in the morning.}
\end{quote}

\emph{Il pescivendolo} — the fishmonger.

\begin{cousinweb}
\emph{pescheria} is built on \emph{pesce} 'fish', which is inherited from Latin \emph{piscis} (EtymDB).
\end{cousinweb}

\begin{grammarlens}[title={What you've built}]
One more word for telling, explaining and describing at B1.
\end{grammarlens}

\subsection*{Your turn}
\begin{itemize}
\raggedright
  \item \emph{Say it:} \emph{la pescheria}
  \item \emph{Say it:} \emph{la pescheria}, once more
  \item \emph{Say it:} a sentence of your own with \emph{la pescheria}
  \item \emph{Recall:} say the Italian for the jeweller's, then the Italian for the delicatessen, then say \emph{la pescheria} again
\end{itemize}

\begin{tcolorbox}[breakable,colback=teal!4,colframe=teal!35!black,title={Before you move on}]
What does la pescheria mean? (\textquotedblleft{}The fishmonger's\textquotedblright{}.) Say it once more.
\end{tcolorbox}
\section[il fruttivendolo]{il fruttivendolo — the greengrocer}

\label{lesson:IT-C454-fruttivendolo}



\begin{quote}
Before the new one: say the Italian for the delicatessen, then the Italian for the fishmonger's.
\end{quote}

\subsection*{You'll want to know: il fruttivendolo}
\textbf{il fruttivendolo} — \textquotedblleft{}the greengrocer\textquotedblright{}.

\begin{quote}
Il fruttivendolo ha le fragole buone.
\end{quote}

\begin{quote}
\emph{The greengrocer has good strawberries.}
\end{quote}

The person and the shop: \emph{vado dal fruttivendolo}.

\begin{cousinweb}
\emph{fruttivendolo} is built from \emph{frutta} 'fruit', inherited from Late Latin \emph{fructa}, and \emph{-vendolo}, from \emph{vendere} 'to sell', inherited from Latin \emph{vēndere} (EtymDB).
\end{cousinweb}

\begin{grammarlens}[title={What you've built}]
One more word for telling, explaining and describing at B1.
\end{grammarlens}

\subsection*{Your turn}
\begin{itemize}
\raggedright
  \item \emph{Say it:} \emph{il fruttivendolo}
  \item \emph{Say it:} \emph{il fruttivendolo}, once more
  \item \emph{Say it:} a sentence of your own with \emph{il fruttivendolo}
  \item \emph{Recall:} say the Italian for the delicatessen, then the Italian for the fishmonger's, then say \emph{il fruttivendolo} again
\end{itemize}

\begin{tcolorbox}[breakable,colback=teal!4,colframe=teal!35!black,title={Before you move on}]
What does il fruttivendolo mean? (\textquotedblleft{}The greengrocer\textquotedblright{}.) Say it once more.
\end{tcolorbox}
\section[la pasticceria]{la pasticceria — the cake shop, the pastry shop}

\label{lesson:IT-C454-pasticceria}



\begin{quote}
Before the new one: say the Italian for the fishmonger's, then the Italian for the greengrocer.
\end{quote}

\subsection*{You'll want to know: la pasticceria}
\textbf{la pasticceria} — \textquotedblleft{}the cake shop, the pastry shop\textquotedblright{}.

\begin{quote}
La domenica compriamo le paste in pasticceria.
\end{quote}

\begin{quote}
\emph{On Sundays we buy pastries at the cake shop.}
\end{quote}

\emph{Le paste} here are pastries, not pasta.

\begin{cousinweb}
\emph{pasticceria} is built on \emph{pasticcio} 'pie', from \emph{pasta} 'dough', which is inherited from Late Latin \emph{pasta}, from Greek \emph{pastá} 'barley porridge' (EtymDB).
\end{cousinweb}

\begin{grammarlens}[title={What you've built}]
One more word for telling, explaining and describing at B1.
\end{grammarlens}

\subsection*{Your turn}
\begin{itemize}
\raggedright
  \item \emph{Say it:} \emph{la pasticceria}
  \item \emph{Say it:} \emph{la pasticceria}, once more
  \item \emph{Say it:} a sentence of your own with \emph{la pasticceria}
  \item \emph{Recall:} say the Italian for the fishmonger's, then the Italian for the greengrocer, then say \emph{la pasticceria} again
\end{itemize}

\begin{tcolorbox}[breakable,colback=teal!4,colframe=teal!35!black,title={Before you move on}]
What does la pasticceria mean? (\textquotedblleft{}The cake shop, the pastry shop\textquotedblright{}.) Say it once more.
\end{tcolorbox}
\section[la rosticceria]{la rosticceria — the rotisserie, the hot-food takeaway}

\label{lesson:IT-C454-rosticceria}



\begin{quote}
Before the new one: say the Italian for the greengrocer, then the Italian for the cake shop.
\end{quote}

\subsection*{You'll want to know: la rosticceria}
\textbf{la rosticceria} — \textquotedblleft{}the rotisserie, the hot-food takeaway\textquotedblright{}.

\begin{quote}
Stasera prendiamo il pollo in rosticceria.
\end{quote}

\begin{quote}
\emph{Tonight we get chicken from the rotisserie.}
\end{quote}

Roast chicken, pizza slices and ready meals.

\begin{cousinweb}
\emph{rosticceria} is built on \emph{arrostire} 'to roast', which is usually traced to a Germanic verb that came into Italian as a borrowing.
\end{cousinweb}

\begin{grammarlens}[title={What you've built}]
One more word for telling, explaining and describing at B1.
\end{grammarlens}

\subsection*{Your turn}
\begin{itemize}
\raggedright
  \item \emph{Say it:} \emph{la rosticceria}
  \item \emph{Say it:} \emph{la rosticceria}, once more
  \item \emph{Say it:} a sentence of your own with \emph{la rosticceria}
  \item \emph{Recall:} say the Italian for the greengrocer, then the Italian for the cake shop, then say \emph{la rosticceria} again
\end{itemize}

\begin{tcolorbox}[breakable,colback=teal!4,colframe=teal!35!black,title={Before you move on}]
What does la rosticceria mean? (\textquotedblleft{}The rotisserie, the hot-food takeaway\textquotedblright{}.) Say it once more.
\end{tcolorbox}
